\chapter{Rudiments de théorie des modèles}

\section{Structures et théories}

\subsection{Langages et structures}

\todo[inline]{blabla}

L'objectif n'étant pas de faire un cours de logique, on se contente de donner
les définitions indispensables pour la suite et de nombreux exemples, le lecteur
pressé peut se contenter des exemples et se référer aux définitions en cas de
doute sur le vocabulaire utilisé. On suit de près le plan de \cite{Marker_2002}
dans lequel on pourra trouver plus de détails.

\begin{defi}[Langage]
Un langage $\cL$ est la donné de:
    \begin{enumerate}[label= (\roman*)]
        \item Un ensemble $\cF$ de fonctions et une arité $n_f \in \bNe$ pour chaque $f \in \cF$.
        \item Un ensemble $\cR$ de relations et une arité $n_R \in \bNe$ pour
        chaque $R \in \cR$.
        \item Un ensemble $\cC$ de constantes.
    \end{enumerate}
Les éléments de ces ensembles sont appelés les symboles de $\cL$.
\end{defi}

\begin{defi}[Structure]
Une $\cL$-structure $\cM$ est la donnée d'un ensemble non vide $M$ appelé
domaine et des interprétations des symboles de $\cL$:
    \begin{enumerate}[label= (\roman*)]
        \item Une fonction $f^{\cM}:M^{n_f} \mapsto M$ pour chaque $f\in \cF$.
        \item Une ensemble $R^{\cM} \subset M^{n_R}$ pour chaque $R \in \cR$.
        \item Un élément $c^{\cM} \in M$ pour chaque $c \in \cC$.
    \end{enumerate}
\end{defi}

\begin{exem}[Langage des groupes]
Si on étudie les groupes, on peut utiliser le langage $\cL_{gp} =
\{\cdot,{}^{-1},e\}$ où $\cdot$ est une fonction binaire ${}^{-1}$ une fonction unaire et $e$ une constante.
Une $\cL_{gp}$-structure $\cG = (G,\cdot^{\cG},{}^{{-1}^{\cG}},e^{\cG})$ sera un ensemble $G$
muni d'une opération binaire $\cdot^{\cG}$, d'une opération unaire
${}^{{-1}^{\cG}}$ et d'un élément distingué $e^{\cG}$. \\
Par exemple $(\bR,\cdot,{}^{-1},1)$ est une $\cL_{gp}$-structure où on interprète
$\cdot$ comme la multiplication usuelle, ${}^{-1}$ comme l'inversion usuelle et
$e$ comme $1$.  \\
Il faut faire attention au fait qu'il n'y a pour le moment aucune règle
sémantique à vérifier, par exemple $\cN = (\bN,+,s,0)$  avec $\cdot^{\cN} = +$ l'addition usuelle, ${}^{{-1}^{\cN}}
= s$ le successeur et $e^{\cG} = 0$ est bien une
$\cL_{gp}$-structure même si ce n'est évidemment par un groupe.
\end{exem}

\begin{defi}[Extensions, Réductions]

\end{defi}

\begin{exem}[Langage des anneaux]

\end{exem}

\subsection{Formules et modèles}

\section{Théories et axiomes}

\section{Ensembles définissables}

\section{Satisfaisabilité et consistence}

\section{Le théorème de compacité}



